\SetPageTheme{theory}
\PageHeader
  {Evaluare FizzBuzz}
  {Aceasta evaluare combina variabile, conditii, repetitie si ordinea regulilor.}
  {FizzBuzz 1 din 4}

\begin{missionbox}[Varianta clasica]
\begin{lstlisting}[style=codequest]
pentru i de la 1 la 30 cu pas 1 executa
    daca (i este multiplu de 3 si 5) atunci
        afiseaza "FizzBuzz"
    altfel daca (i este multiplu de 3) atunci
        afiseaza "Fizz"
    altfel daca (i este multiplu de 5) atunci
        afiseaza "Buzz"
    altfel
        afiseaza i
    sfarsit
sfarsit
\end{lstlisting}
\end{missionbox}

\begin{conceptbox}[Ordinea regulilor]
Cazul ``si 3, si 5'' trebuie verificat primul. Altfel, valorile comune sunt capturate prea devreme de una dintre regulile simple.
\end{conceptbox}

\clearpage

\SetPageTheme{guided}
\PageHeader
  {Urmarire pe cazuri}
  {Testam explicit valori cheie pentru a verifica ordinea corecta.}
  {FizzBuzz 2 din 4}

\begin{guidebox}[Tabel de control]
Completeaza:
\begin{center}
\begin{tabular}{llll}
\toprule
Valoare & Regula activa & Output & De ce \\
\midrule
3 & & & \\
5 & & & \\
15 & & & \\
18 & & & \\
20 & & & \\
30 & & & \\
\bottomrule
\end{tabular}
\end{center}
\AnswerSpace{6}
\end{guidebox}

\clearpage

\SetPageTheme{challenge}
\PageHeader
  {Scor si variatii}
  {Acum adaugi punctaj si reguli noi fara sa pierzi claritatea.}
  {FizzBuzz 3 din 4}

\begin{solobox}[FizzBuzz cu punctaj]
Construieste varianta:
\begin{itemize}
\item \texttt{Fizz} = 1 punct;
\item \texttt{Buzz} = 1 punct;
\item \texttt{FizzBuzz} = 3 puncte.
\end{itemize}
La final afiseaza scorul total.
\AnswerSpace{8}
\end{solobox}

\begin{solobox}[Ping/Pong]
Propune regula:
\begin{itemize}
\item \texttt{Ping} pentru multiplii lui 4;
\item \texttt{Pong} pentru multiplii lui 6;
\item \texttt{PingPong} pentru multiplii ambilor.
\end{itemize}
\AnswerSpace{7}
\end{solobox}

\clearpage

\SetPageTheme{challenge}
\PageHeader
  {Pagina de rationament}
  {Scrii explicatia regulilor, nu doar codul final.}
  {FizzBuzz 4 din 4}

\begin{solobox}[Explica ordinea]
In 6-8 randuri, explica:
\begin{enumerate}
\item de ce exista un caz special combinat;
\item de ce ordinea verificarilor poate strica rezultatul;
\item cum ai testa rapid daca varianta ta este corecta.
\end{enumerate}
\AnswerSpace{12}
\end{solobox}

\clearpage
